\subsubsection{Conclusion}

Cette manipulation a permis de mesurer la vitesse du son dans l'air et dans le CO$_2$ avec le trombone de Koenig. Pour les deux gaz, $\Delta x$ est proportionnel à $p/(2f)$, la vitesse du son est donc, qualitativement, indépendante de la fréquence dans les domaines étudiés ($1\,500$ à $3\,000$ Hz pour l'air, $1\,000$ à $2\,000$ Hz pour le CO$_2$).

\begin{table}[H]
    \centering
    \small
    \begin{tabular}{lccccc}
        \toprule
        Gaz & $T$ [K] & $v_{\text{exp}}$ [m/s] & $v_{\text{th}}$ [m/s] & Écart [m/s] (\%) & Incert. combinée [m/s] \\
        \midrule
        Air    & 297,1 & $335{,}3 \pm 2{,}0$ & $345{,}3 \pm 0{,}9$ & 10,0 (2,9) & 2,2 \\
        CO$_2$ & 296,3 & $272{,}2 \pm 1{,}5$ & $268{,}3 \pm 0{,}7$ & 3,9 (1,5)  & 1,7 \\
        \bottomrule
    \end{tabular}
    \caption{Tableau récapitulatif, conclusion, manipulation 2}
    \label{tab:conclusion-trombone}
\end{table}

Le \refit{Tableau}{\ref}{tab:conclusion-trombone} résume les résultats pour les deux gaz. Dans les deux cas, l'écart est supérieur à l'incertitude combinée et la \refit{loi}{\eqref}{eq:vitesse-theorique} est validée aux incertitudes expérimentales. Pour le CO$_2$, les causes supposées de l'écart sont une présence résiduelle d'air dans le tube, qui augmenterait la vitesse mesurée, et le petit nombre de maxima ($p = 2$ ou $3$).

L'objectif est donc atteint, les vitesses sont mesurées avec une incertitude relative de $0{,}6\,\%$ et le modèle est confirmé, aussi bien qualitativement (indépendance de la vitesse vis-à-vis de la fréquence, proportionnalité de $\Delta x$ à $p/(2f)$) que quantitativement, avec des écarts de $2{,}9\,\%$ (air) et $1{,}5\,\%$ (CO$_2$) par rapport aux valeurs théoriques. Les mesures dans l'hélium n'ont pas été réalisées car le gaz n'était pas disponible.
